% Abstract
% Standalone, compilable draft (FPT-template article class + booktabs), sibling to
% sec_1_introduction.tex / sec_9_conclusion.tex. One dense paragraph (~200-250 words)
% compressing the whole contribution story: the modular near-annotation-free
% detection+completion pipeline, the detection headline and its root-caused recall
% limit, the synthetic-pretraining-redundancy negative, and the completion half's
% scientific core (Chamfer-invalidity + the leakage-free donor metric + the
% PARTIALLY-HOLDS held-out result). It is the first thing a committee reads, so every
% number and hedge must be defensible.
%
% MERGE/NUMBERING NOTE: this file opens with \section*{Abstract} (unnumbered) for
% standalone compilation; at final assembly the assembler places it in the front
% matter, before Chapter 1. It cites no chapters and no literature -- an abstract is
% self-contained prose -- so there are no cross-references or \cite{} here.
%
% Honesty-phrasing rules applied (personal_agenda P1):
%   - recall is the car-only, >=10-surviving-point figure (Sec 4.1); the abstract does
%     NOT imply recall against all annotated cars (~0.54, Finding #48)
%   - the held-out completion result is PARTIALLY HOLDS, never "confirmed"/"holds"
%   - the detection headline is a single-sequence (seq 08) result, stated as such
%   - negatives (recall repairs, synthetic redundancy) are named as findings
%   - the donor metric is "a validated, leakage-free" metric, NOT "the first"
%     (no literature-priority overclaim; matches C7 and the Ch 9 review)
%   - claims kept within the Section-6 claims map (C1-C11); no mover accuracy claim
%
\chapter*{Abstract}
\label{ch:abstract}

This thesis studies how far an interpretable, near-annotation-free pipeline can go
at detecting and completing cars in LiDAR point clouds, using classical geometry
where geometry is reliable and small learned models only where discrimination
needs it: ground removal and HDBSCAN density clustering propose objects,
a lightweight binary classifier labels cars, a tracker filters flicker, and a PCN
network completes each detected car's occluded surface. On sequence~08 (the reported
operating point, \num{4071} scans) the detector reaches precision \num{0.905},
recall \num{0.730}, and $F_1$ \num{0.808} at a point-IoU of \num{0.3}, with recall
measured against the car-only, at-least-ten-surviving-points denominator;
leave-one-sequence-out cross-validation over all eleven labelled sequences confirms
this is not a model-selection artefact (leakage-free seq-08 recall \num{0.737},
pooled recall \num{0.719} across the eleven) and reports the per-sequence spread.
The recall
shortfall is root-caused, not tuned away: density clustering splits single cars into
fragments (\SIrange{31}{37}{\percent} of ground-truth cars), and a search of repair
strategies fails to generalise, reported as a chain of negative findings; synthetic
classifier pretraining is likewise shown redundant once real training clusters are
plentiful, the sim-to-real gap being total and symmetric. On the completion side,
reference-based Chamfer distance against an accumulated pseudo-ground-truth is shown
invalid on real cars---the raw partial input scores best---so the thesis defines and
validates a leakage-free donor-frame occluded-side coverage metric. Through it,
completion recovers genuinely unseen surface several times above a symmetry-mirror
baseline (occluded-side coverage \num{0.000}~$\rightarrow$~\num{0.304}) and improves
amodal boxes on static, gate-passed cars, a benefit that partially holds under a
pre-registered held-out test. The contributions are methodological and diagnostic,
not architectural, and each is defended by explanation and pre-registration rather
than a further point of metric improvement.
